\begin{table}[htbp]
\centering
\small
\caption{V1 E1（表現デコードピークの局在）：各モデルファミリーにおける感情認識（Reader）および自己報告（Self）の最高線形デコード性能（$R^2$、Pearson $r$）と最適相対深度 $d^*$。}
\label{tab:v1_peak_decodability}
\begin{tabular}{lll cccc c}
\toprule
 & & & \multicolumn{2}{c}{\textbf{Valence Peak}} & \multicolumn{2}{c}{\textbf{Arousal Peak}} & \textbf{Reader--Self} \\
\cmidrule(lr){4-5} \cmidrule(lr){6-7}
\textbf{Family} & \textbf{Variant} & \textbf{Task} & Layer ($d^*$) & $R^2$ ($r$) & Layer ($d^*$) & $R^2$ ($r$) & Distance $\Delta d^*$ \\
\midrule
\multirow{4}{*}{\textbf{Qwen 2.5 1.5B}} & \multirow{2}{*}{Base} & Reader & L22 (0.81) & 0.378 (0.617) & L24 (0.89) & 0.130 (0.369) & \multirow{2}{*}{V:0.00, A:0.07} \\
                                 &  & Self   & L22 (0.81) & 0.391 (0.626) & L26 (0.96) & 0.123 (0.356) &  \\
\cmidrule(lr){2-8}
                                 & \multirow{2}{*}{Instruct} & Reader & L15 (0.56) & 0.215 (0.467) & L17 (0.63) & 0.073 (0.301) & \multirow{2}{*}{V:0.07, A:0.11} \\
                                 &  & Self   & L13 (0.48) & 0.285 (0.536) & L20 (0.74) & 0.075 (0.300) &  \\
\midrule
\multirow{4}{*}{\textbf{Llama 3.2 1B}} & \multirow{2}{*}{Base} & Reader & L14 (0.93) & 0.239 (0.490) & L4 (0.27) & 0.068 (0.288) & \multirow{2}{*}{V:0.20, A:0.00} \\
                                 &  & Self   & L11 (0.73) & 0.243 (0.494) & L4 (0.27) & 0.065 (0.282) &  \\
\cmidrule(lr){2-8}
                                 & \multirow{2}{*}{Instruct} & Reader & L6 (0.40) & 0.317 (0.564) & L8 (0.53) & 0.113 (0.352) & \multirow{2}{*}{V:0.13, A:0.00} \\
                                 &  & Self   & L8 (0.53) & 0.355 (0.597) & L8 (0.53) & 0.106 (0.346) &  \\
\midrule
\multirow{4}{*}{\textbf{Gemma 3 1B}} & \multirow{2}{*}{Base} & Reader & L16 (0.64) & 0.331 (0.576) & L15 (0.60) & 0.043 (0.251) & \multirow{2}{*}{V:0.00, A:0.04} \\
                                 &  & Self   & L16 (0.64) & 0.330 (0.576) & L16 (0.64) & 0.044 (0.254) &  \\
\cmidrule(lr){2-8}
                                 & \multirow{2}{*}{Instruct} & Reader & L12 (0.48) & 0.301 (0.552) & L12 (0.48) & 0.094 (0.323) & \multirow{2}{*}{V:0.08, A:0.36} \\
                                 &  & Self   & L14 (0.56) & 0.375 (0.613) & L21 (0.84) & 0.088 (0.321) &  \\
\midrule
\multirow{4}{*}{\textbf{OLMo 2 1B}} & \multirow{2}{*}{Base} & Reader & L14 (0.93) & 0.418 (0.647) & L7 (0.47) & 0.089 (0.314) & \multirow{2}{*}{V:0.00, A:0.13} \\
                                 &  & Self   & L14 (0.93) & 0.458 (0.677) & L9 (0.60) & 0.106 (0.338) &  \\
\cmidrule(lr){2-8}
                                 & \multirow{2}{*}{Instruct} & Reader & L7 (0.47) & 0.327 (0.573) & L11 (0.73) & 0.085 (0.307) & \multirow{2}{*}{V:0.00, A:0.27} \\
                                 &  & Self   & L7 (0.47) & 0.352 (0.594) & L7 (0.47) & 0.086 (0.310) &  \\
\bottomrule
\end{tabular}
\vspace{1ex}
\begin{minipage}{\linewidth}
\footnotesize
\textbf{Note:} ReaderとSelfのdecodability peakは多くの条件で近接したが、一部のmodel / axisでは相対深度0.2以上の乖離も観測された。したがって両taskはaffect-relevant informationを共有する一方、その最適なreadout depthは完全には一致しない。
\end{minipage}
\end{table}